\documentclass[12pt,a4paper,onecolumn]{article}

% Encoding and language
\usepackage[utf8]{inputenc}
\usepackage[T1]{fontenc}
\usepackage[english]{babel}

% Font
\usepackage{times}

% Set margins
\usepackage[a4paper, top=2cm, bottom=2cm, left=2.5cm, right=2.5cm]{geometry}

% Tabulation
\usepackage{tabto}

% Line spacing
\usepackage{setspace}
\onehalfspacing

% Page numbers
\usepackage{fancyhdr}
\pagestyle{fancy}
\fancyhf{} % Clear default header/footer
\rfoot{\thepage} % Right-aligned page numbers in footer

% Remove page number from cover/abstract/etc
\usepackage{titlesec}
\newcommand{\noPageNumber}{\thispagestyle{empty}}

% Prevent indentation
\setlength{\parindent}{0pt}
\setlength{\parskip}{0.5em}

% Figure management
\usepackage{graphicx}
\usepackage{caption}
\usepackage{subcaption}

\usepackage{pgfplots}
\pgfplotsset{compat=1.16}
\def\mathdefault#1{#1}
\everymath=\expandafter{\the\everymath\displaystyle}

% Subfigure management
\usepackage{layouts}
%\usepackage{floatrow}
%\usepackage[label font=bf,labelformat=simple]{subfig}% <-- changed
%\floatsetup[figure]{style=plain,subcapbesideposition=top}


% Optional: headings formatting
\titleformat{\section}{\large\bfseries}{\thesection}{1em}{}
\titleformat{\subsection}{\normalsize\bfseries}{\thesubsection}{1em}{}

% Optional: remove widow/orphan lines
\widowpenalty=10000
\clubpenalty=10000

% Mathematics
\usepackage{amsmath,mathtools,amssymb} 

\begin{document}
 %% textwidth in cm: \printinunitsof{in}\prntlen{\textwidth}

% Cover page (no page number)
\noPageNumber
\begin{titlepage}
\centering
\vspace*{3cm}
{\Huge Evaluating the impact of maturation time variability in stage-structured host-parasitoid models (cette page de garde va changer) \par}
\vspace{2cm}
{\Large Mathis Drieu \par}
\vfill
February-June 2025
\end{titlepage}

% Abstract (no page number)
\noPageNumber
\begin{abstract}
xxx
\\
xxx
\\
\\
xxx
\\ 

\end{abstract}

% Table of Contents, List of Abbreviations (no page numbers)
\noPageNumber
\tableofcontents


% \listoffigures % If needed
% \listoftables  % If needed

% List of abbreviations
\section*{List of Abbreviations}
\noPageNumber
% Your content

% Main content starts here (with page numbers from 1)
\newpage
\setcounter{page}{1}






\section{Introduction}

Je suis pas très content de cette intro je balance beaucoup de trucs un peu douteux et c'est mal formulé. Mais j'ai essayé d'appliquer les conseils de Quentin à savoir : plus repartir d'un enjeu générale en écologie, pour arriver en entonnoir à une pbtmq précise (de base j'avais fait un sorte de revue qui allait un peu nul part). \\ 

Donc surtout les premiers paragraphes retenez plutôt l'intention que le fond du propos... \\

\tabto{1 cm}[Agricultural pest species [Insects] constitute a fundamental challenge [economic loss?] facing food production in all agroecosystems]. Biocontrol is a field of pest control that relies on biological agents to curb pest populations. This field is [old but quickly developing] notably in agroecology, one of its main goals being to avoid, or at least reduce, the use of industrial insecticides. One of the primary levers of biocontrol is the introduction of insect parasitoids that target pest species. They induce the death of their host by laying eggs [inside them : reformuler] at various stages of their developmental cycle. The host's body provides shelter and food for layed parasitoid juveniles, that eventually emerge and reproduce thus [closing the cycle : reformuler]. Mathematical models have been developed to model the evolution of two interacting populations : the pest and its enemy. More specifically, in the case of host-parasitoid interactions, the parasitoid constitutes a supplementary mortality factor for the host and is dependent on the host population for its own reproduction. Regarding the continuous property of parasitism [et le fait que le modèle discret de Nicholson Bailey est instable???], continuous time models have been widely used in host-parasitoid modelization and [have relied on the set of a term giving the rate of parasitism through time as a function of parasitoid population] [source? nulle la phrase non?]

\tabto{1 cm} [In any case : reformuler], both host and parasitoid are insect species with discontinuous development. Successive developmental stages are not identically sensitive to the various ecological interactions of the species. In the case of parasitism, it can occur on one or several specific stages of the host (Murdoch et al. 1987, ? ). Hence, stage-structured (or class-structured models) models have been increasingly used to aggregate individuals having certain characteristics, for example, into age classes or size classes, and that critically have the same set of interactions with their environment.  [tout ça manque de source mais Nisbet et Gurney 1983?]. The set of stages raises the challenge of modeling the transition between them, notably the timing of the transition.

\tabto{1 cm}The easiest way to account for the transition between stages is to write that the rate of change of the stage population $x$ is proportional to the stage population.  In other words, $\frac{dx}{dt} = -mx$, m is the maturation rate of the stage and 1/m is the mean maturation time. This framework allows parametrizing the timing of transition but it does not allow to control the variability of maturation time inside the stage population. The ODE framework implicitly assumes an underlying stochastic state transmission model with an exponentially distributed dwell time in the stage [Hurtado]. Particularly, it assumes that the probability of transition is maximal at the entrance of the stage. Robertson and co-workers (2018) showed that about half of the stage-structured models found in recent literature use this Ordinary Differential Equations (ODE) framework to account for the transition between states. In line with the work of Gurney and co-workers in 1983, the remaining models by Robertson et al. (2018) majorly used fixed delays (DDE). These models assume that all the individuals of a given cohort mature synchronically, i.e. that the maturation delays are fixed. This means that the models do not account for within-species (within-stage) variability in maturation time.  Overall, despite their flexibility on the parametrization of mean maturation time, the two approaches (ODE and DDE)  rely on fundamentally different assumptions on maturation time variability. However, the latter can greatly influence the models' predictions. 

\tabto{1 cm} The use of fixed delays was shown to be a destabilizing factor of community dynamics as soon as the 1930s (Bailey 1931, cited in Smith and Mead 1974). Seeking higher biological realism, developmental time variability was subsequently introduced and shown to be an important stabilizing factor of community dynamics, including pest populations (Smith and Mead 1974, Blythe et al 1984, Bellows 1986, Briggs 1993, Abrams and Quince 2005, Xu 2010, de Valpine 2014, Robertson 2018), and notably so with experimental support (Cronin et al 2016). In the case of host-parasitoid models, this variability was considered a form of parasitism heterogeneity (or risk aggregation) (Xu et al. 2010). This concept is characterized by [represents?] the variance [variability] of parasitism risk among hosts. [Ici on peut introduire : "on peut créer ce type de variance en faisant l'hypothèse d'une distribution spatiale des hôtes en patches. Le terme k, dit d'overdispersion, représente alors l'homogénéité du risque de parasitisme... cf May 1978 qui l'explique mieux. Ca peut être utile pour jusitifer le fait que plus tard on regarde l'équivalence k-n]. Parasitism heterogeneity was shown to stabilize populations (Chesson and Murdoch 1984, Hassel 2001ab, Taylor 1993)

\tabto{1 cm} However, Wearing et al. 2004 showed that their host-parasitoid model' predictions based on the fixed-development approach are differentially sensitive to variability and in key life stages of insects. The stage to which developmental time variability is added, the studied dynamic, and the used set of parameters (and their distance to a pre-existing bifurcation) were shown to influence the effect of modified developmental variability. These results suggest that the effect of developmental variability should be studied not only on stability but also on persistence, as well as population equilibria, general properties of oscillations when they exist, etc. 

\tabto{1 cm} Besides, models' sensitivity to maturation variability stresses the need to estimate empirical developmental variability beforehand to properly model given host-parasitoid dynamics. A clear review of the methods used to assess developmental variability was given by de Valpine and co-workers (2014). Empirical distributions of maturation time often display a delay where no maturation occurs, a clear mode as well as a relatively pronounced right-skewness. 
\tabto{1 cm} 
[Hence, It can be argued that] [On pourrait dire que... Avec cette idée de spectre, j'essaie ici de retrouver cette idée d'axe de variabilité développementale à explorer facilement avec le n du LCT. J'essaie donc d'introduire les bords de l'axe que je vais explorer avec la LCT. Mais je garde bien à l'esprit que c'est une manière parmi d'autres d'explorer cette variabilité qui ne se limite pas à cet axe : la loi Erlang est arbitraire, on ne la maximise jamais, n=1 est un cas particulier très arbitraire aussi]. Realistic maturation time distributions [therefore] lie at various positions of a variability spectrum which extrema can be represented by the ODE framework on one side (high variability) and the DDE framework on the other (no variability).  The aim of this study is therefore : 1. To assess the position of various pest and parasitoid species on this maturation variability spectrum ; 2. To elucidate how the position on this spectrum influences the predictions made by host-parasitoid models. 

\tabto{1 cm} To explore this spectrum, arbitrary maturation time distributions can be used in integrodifferential equations (IDEs). These systems quickly lead away from analytical methods towards computational ones (de Valpine 2014). Hence, modelers have relied on a few main distributions. Gamma (or Erlang, a special case where the shape parameter is a positive integer) and Normal distributions have been widely used in distributed development models of insect populations (SOURCE? deValpine, Robertson?, Wearing?) [Enoncer l'existence des degré jours et des approches de ce style?]. They allow for a sufficiently realistic hypothesis about stage transitions with relatively few parameters associated. It is however possible to account for more realistic stage transitions while also constructing models that allow for relatively simple algebraic analysis. Erlang distributions can implicitly be obtained by dividing each stage into n sub-stages governed by a system of n coupled ODEs (MacDonald 1978, Hurtado 2018). This is called the Linear Chain Trick (LCT) and has been increasingly used in ecological (Smith and Mead 1974, McNair 1995, Wearing 2004, Abrams and Quince 2005, Bjornstad et al. 2016, Nelson et al. 2019) and epidemiological studies (Wearing 2005, Hurtado 2019).
 [More generally, Hurtado and Kirosingh (2019) suggested a method to bypass the IDE framework and allow modelers to assume maturation (or "passage") time distributions from a family called the phase-type distributions. These distributions can be thought of as the family of all possible hitting time distributions for continuous-time Markov chains. In the case of a "linear chain" of ODEs, the Erlang distribution is involved. As suggested by the example SEIR model with heterogeneity among infected individuals shown by Hurtado and Cameron (2021), this new framework allows highly flexible ODE models assuming more realistic maturation time distribution and can be computationally more efficient. Even outside this framework, the LCT can be judicious. Nelson et al. (2019) used the LCT to implement hierarchical asymmetric larval competition : they assumed the per-capita mortality rate to depend only on the density of animals that are larger or equal to the size of each individual. The competition mortality terms uniquely depend on the downstream LCT subcompartments. -> tout cet encadré : plutôt en discussion].
\\ 
D'ailleurs ces deux paragraphes précédents pourraientt être complétés par une figure qui montre justement l'axe=le spectre de variab dev, une version plus propre de la figure ci-dessous
\setcounter{figure}{-1}
\begin{figure}[!ht]
\centering	
\includegraphics[width=0.8\textwidth]{/home/mathisdr/Desktop/2024-2025/Stage/Internship_2025_LCT/model_host_parasitoid/figs/pbmtq_figure.png}
\caption{Exploring the effect of developmental (time) variability on host-parasitoid models using the linear chain trick.}
\label{fig:intro}
\end{figure}




\tabto{1 cm} Murdoch et al. (1987) and Godfray and Hassel (1989) constitute two of the first relatively simple models that succeeded in predicting (qualitative) dynamics observed in natural host-parasitoid systems. Their simplicity allowted to greatly contribute to the understanding of the complex dynamics emerging from interactions between hosts and parasitoids [général, gratuit, pas sourcé :( ]. These models were built with fixed maturation delays with the justification, introduced by Blythe et al. (1984), that "the witdh of the distribution [...] is short in comparison with the mean developmental time" [je laisse ça? on s'en fout? comment je tourne ça ? Parce que en gros ils le justifient en mode c bon blythe a dit ça passe mais donc peut-être que pr les espèces qui les intéressent (notamment la red scale pour Murdoch) c'est ok mais c'est pas généralisable. Ou peut-être même qu'ils ont jamais eu vraiment de scrupule et que c'est un peu un argument d'autorité]. Thus, the first part of this study aims to evaluate the realism of this hypothesis. We will review the number of subcompartment $n$, i.e. the Gamma law shape, that corresponds to maturation time variability in studies that experimentally measure insects' development times. In the second part, we will apply the LCT to these models, yielding distributed-delay equivalents of them. We will use the number of subcompartments, n, to easily explore the maturation time variability spectrum, thereby assessing the robustness of the models' prediction to manipulation of this new parameter.

\section{Models and methods}
\subsection{Survey of Erlang/Gamma distribution shape (n)}
\tabto{1 cm} Wagner et al. (1984) introduced a method to fit a Weibull distribution to normalized maturation time of individually marked insects bred in laboratory at various temperatures. As times are normalized (divided) by the median maturation time at the given temperature, the distribution is built to only represent the variability of maturation time (independently of its mean). Xu et al. (2010) reviewed the literature using Wagner's method. They gathered the Weibull parameters of various insects' developmental stages, and fitted a shifted Gamma distribution. They simulated 5000 time variates from the given Weibull laws and found the Gamma parameters that maximized the likelihood of their simulated data. They found that the estimated shapes (n) were either relatively small (1-10) or larger (>20). Here, the same method was applied but according to two constraints on the shift of the fitted Gamma law : it was estimated either in the range [0;1] (similarly to Xu et al.) or forced to 0. Contrary to Xu et al, combined stages were not eliminated. Overall, around 20 papers were added to the review data set originally made up of around 20 papers in Xu et al. Similarly to Xu et al., the majority of the estimates were obtained on pest species of some economic importance as well as on their natural enemies. The dataset involves one parasitoid : Cotesia Melanoscela (Gould and Elkinton 1990). Tables and Python (Python 3.12.2) script summarizing the survey can be found on [GITHUB?]. Visual representations of the simulated Weibull variates and the corresponding fitted Gamma distributions can be found in the file. Based on these figures, the goodness-of-fit was deemed generally sufficient.

\subsection{Murdoch (1987) model}
[Changement à faire probable : moins de théorie sur ce qu'a trouvé Murdoch (à remettre dans les résultats, plus d'explication de stability diagram etc]
\tabto{1 cm} We consider the continuous-time host-parasitoid model detailed in Murdoch et al. 1987 in a DDE approach. The host and parasitoid life cycles are divided into 2 stages : juvenile ($H_1, P_1$) and reproductive adult ($H_2, P_2$). Parasitism occurs on the juvenile host stage $H_1$ under a random search hypothesis. The application of the LCT to obtain a distributed-delay model equivalent of the fixed-delay model introduced by Murdoch is detailed in Appendix A1, as well as the non-trivial equilibrium densities. From now on, except in Figure 2c, the temporal scale is normalized by the maturation time of the host : $m_H = 1/T_{H1} = 1$.
\\
FIGURE : COMPARTIMENTS, SYSTEME MATH LCT EQUATIONS
\\
\tabto{1 cm} Local stability was numerically estimated using the function Eigenvalues of Wolfram Mathematica (Version 14.1) applied to the Jacobian matrix of the LCT system. Numerical simulation was realized using the Mathematica function NDSolve with an automatic (and dynamically changing) solver method. In the light of the frequent stiffness issues of the simulations, NDSolve was found to offer the most fast and reliable solver compared to other tested methods in Python. For further details, relevant Mathematica and Python scripts can be found on [GITHUB?].

\subsection{Godfray and Hassel (1989) model}
\tabto{1 cm} Godfray and Hassel (1989) introduced a slightly complexified version of the Murdoch 1987 DDE model. First, two invulnerable juvenile stages ($H_1$ and $H_3$) were added for the host and surround the vulnerable stage $H_2$. Second, the parasitism term was changed to involve density-dependence to $P_2$. Parasitism thus saturates with the (local) density of adult parasitoids : this is called pseudo-interference. The LCT model equivalent to the DDE model is detailed in [Appendix 2]. It was built and analyzed analogously to the Murdoch 1987 LCT model using the Jacobian of the LCT system (see [GITHUB XXX]). From now on, Erlang shape parameters are respectively $n_1, n_2, n_3, m$ for $H_1, H_2, H_3, P_1$ maturation times.

\section{Results}
\subsection{Survey of Erlang/Gamma distribution shape (n)}

\begin{figure}[!ht]
\centering
	\begin{subfigure}[t]{0.48\textwidth}
	  \centering
	  \includegraphics[width=\textwidth]{/home/mathisdr/Desktop/2024-2025/Stage/Internship_2025_LCT/survey_gamma_shape/figs/n_survey_not_clipped.pdf}
	  \caption{}
	  \label{fig:a}
	\end{subfigure}
	\hfill
	\hfill
	\begin{subfigure}[t]{0.48\textwidth}
	  \centering
	  \includegraphics[width=\textwidth]{/home/mathisdr/Desktop/2024-2025/Stage/Internship_2025_LCT/survey_gamma_shape/figs/n_survey_clipped.pdf}
	  \caption{}
	  \label{fig:b}
	\end{subfigure}

\vspace{1em} % Optional: vertical space between rows

	\begin{subfigure}[t]{0.5\textwidth}
	  \centering
	  \includegraphics[width=\textwidth]{/home/mathisdr/Desktop/2024-2025/Stage/Internship_2025_LCT/survey_gamma_shape/figs/n_survey_example_species.pdf}
	  \caption{}
	  \label{fig:c}
	\end{subfigure}


\caption{Survey of Gamma shape parameter of insects stages found in literature. In \textbf{(a)} and  \textbf{(b)}, only stages like Egg, Larva, Pupa/Nymph and Proto/Deutonymph were retained. \textbf{(a)} : Fitted gamma shape parameters (n). Shift was either constrained to [0;1] (top) or to 0 (bottom). \textbf{(b)} : Fitted n with all n>100 counted as 100. \textbf{(c)} : Fitted n for stages or whole juvenile development of a few insects involved in agroecosystems. Ips avulsus and Feltiella acarisuga have only 1 nymphal stage. Grey bars indicate that the Gamma distribution of the corresponding stage was not fitted due to a negative shift in the Weibull law estimates given in the original paper.}
\label{fig:Survey}
\end{figure}

\tabto{1 cm} The average Erlang shape (n) obtained is 96.0 if the fitted shift is constrained in $[0;1]$ and 123.3 if it is forced to 0 (Figure 1a). The minimal shapes are respectively 1.9 and 4.5. The maximum shape is 782.2 regardless of the constraint on the shift. These ranges are substantially greater than the one indicated in Xu et al. (2010). The lines connecting the two violins in Figure 1a;1b indicate that the constraint to have a null shift is associated with a higher fitted n. [SUPP FIGURE 1?] gives a representative example of such a case. For 20 out of 57 fitted stages (35\% of cases), the difference between the fitted n with constraint $\tau = 0$ and the fitted n with constraint $\tau \in [0;1]$ is superior to 1. 
\tabto{1 cm} n display a large variability (Figure 1a) and appears to be dependent on the studied species, as suggested by the 5 species detailed in Figure 1c. [QUANTIFIER LA VARIABILITE INTRAESPECE vs INTERESPECE?]. Figure 1c also suggests that maturation variability is generally lower when considering combined stages such as the complete development of insects [MIEUX QUANTIFIER CA?]. 

\subsection{Murdoch (1987) model}
\tabto{1 cm} Using fixed delays, W.W. Murdoch showed that the invulnerability of the host adult reproductive stage stabilizes the system because this this stage constitutes a "refuge" effect. In line [reformuler : de manière cohérente], the longevity of the adult host was found to be the main stability factor of the system.  Here, maturation variability was added using the LCT.  [transi abrupte avec la suite]
\tabto{1 cm} Maturation variability causes the non-trivial equilibrium densities to increase, except for the case of $H_1$ (Figure 2a). The increase is particularly sharpe when n is near 1. For instance, between (n,m) = 1 and (n,m) = 5, $H_2^*$ is divided by 5 for this set of parameters. Between (n,m) = 10 and (n,m) = 100, the ratio reduces down to 1.13 . [\textit{D'ailleurs cette explosion des eq à n très bas est encore plus brutale pour Godfray 1989 avec les paramètres de Xu : soit c un effet des paramètres, soit c un effet de la modification du modèle. Si c le modèle qui fait ça, soit c le nombre de stade, soit c le terme de parasitisme. Si c le nombre de stade (ce que je soupçonne car en gros quand on calcule les équilibres NT, on trouve un equilibre pour le stade d'hôte parasité et la hauteur de l'eq dépend du parasitisme. Ensuite on trouve tous les autres stades de H. Notamment, on 'remonte' les stades en amont et on trouve finalement une population de H1 d'autant plus grande qu'il y a de mortalité de H1 à Hparasité typiquement. Donc par exemple si H9 est parasité il me semble que H1* sera très grand? à vous de me dire si je dis des bêtises, et si c'est pertinent de le remarquer en discussion}]. In the ODE framework (n = m = 1), the system is almost always stable, except for a short range of low $T_{H2}/T_{H1}$ values. The unstable area increases as (n,m) increase and does not change much after (n,m) attain values around 20 or 30 (Figure 2b). The bifurcation diagram (Figure 2c) specify the oscillation properties of the system. In addition to increasing the area of unstability, the increase of (n,m) causes the increase of the amplitude and the period of oscillations. This extent of this effect is dependent of the chosen parameters [SUPPL. FIGURE ??]. Finally, Figure 2d suggests that host fecundity is stabilizing for $ n \leq 2 $ and else destabilizing. Altough the parasitoid' maturation variability is stabilizing (SUPPL FIG?), it does not reverse the effect of fecundity on stability (SUPPL. FIGURE avec les courbes pour m=10?). Together, these results indicate that maturation variability is indeed a stabilizing factor and can modulate the effect of other demographic parameters (here fecundity) on stability. 


\begin{figure}[!ht]
\centering
	\begin{subfigure}[t]{0.48\textwidth}
	  \centering
	  \includegraphics[width=\textwidth]{/home/mathisdr/Desktop/2024-2025/Stage/Jupyter_stage/Figures/Murdoch_equilibrium.pdf}
	  \caption{}
	  \label{fig:a}
	\end{subfigure}
	\hfill
	\begin{subfigure}[t]{0.48\textwidth}
	  \centering
	  \includegraphics[width=\textwidth]{/home/mathisdr/Desktop/2024-2025/Stage/Jupyter_stage/Figures/Murdoch_stability_diagram.pdf}
	  \caption{}
	  \label{fig:b}
	\end{subfigure}

\vspace{1em} % Optional: vertical space between rows

	\begin{subfigure}[t]{0.48\textwidth}
	  \centering
	  \includegraphics[width=\textwidth]{/home/mathisdr/Desktop/2024-2025/Stage/Jupyter_stage/Figures/Murdoch_bifurcation_diagram_Wearing-par.pdf}
	  \caption{}
	  \label{fig:c}
	\end{subfigure}
	\hfill
	\begin{subfigure}[t]{0.48\textwidth}
	  \centering
	  \includegraphics[width=\textwidth]{/home/mathisdr/Desktop/2024-2025/Stage/Jupyter_stage/Figures/Murdoch_fecundity.pdf}
	  \caption{}
	  \label{fig:d}
	\end{subfigure}

\caption{Effect of maturation variability on different properties of the Murdoch 1987 model. \textbf{(a)} : Non trivial equilibrium densities as a function of the Erlang distribution shapes n and m. Parameters are $\delta_{H1} = 0, \delta_{P1} = 1,  \delta_{P2} = 8, \rho = 33, a = 1, m_P = 1, \delta_{H2} = 1/T_{H2} = 0.2$. The system displays a limit cycle to the right of the dashed red line (for n,m > 10).  \textbf{(b)} : Local stability diagram of the system around the non-trivial equilibrium. The system is unstable to the left of the curves, and stable to their right. \textbf{(c)} : Bifurcation diagram of the system with parameters inspired from those in Wearing 2004 : $\delta_{H1} = 0.00425, \delta_{P1} = 0.1, \delta_{P2} = 0.5, \rho = 105, m_P = 0.05,  T_{H2} = 1/\delta_{H2} = 5, T_{H1} = 1/m_H = 37$. \textbf{(d)} : Local stability diagram of the system around the non-trivial equilibrium showing the effect of the average number of offspring of an adult host $\rho$.}
\label{fig:Murdoch}
\end{figure}

\subsection{Godfray and Hassel (1989) model}
\tabto{1 cm} The adjunction of a non-linear term by Godfray and Hassel (1989) on a Murdoch-type model caused the system to display a new type of oscillations which period is roughly the same as the generation time of the host. These (single-) generation cycles will hereafter be named "GC". Their model predict that GCs are possible when the ratio between the parasitoid generation time and the host generation time (hereafter named "ratioT") is non-integer (0.5, 1.5, ...).  Mathematically, this is caused by areas delimited by new stability curves arising from new couples of conjugated characteristic roots (of the characteristic equation of the system) that can have positive reals parts (Godfray and Hassel 1989, Xu et al. 2010). To study these new areas (i.e the occurrence of GCs) under a distributed delay assumption, the equivalent LCT model was built and the stability curves were studied.
\begin{figure}[!ht]
\centering
	\begin{subfigure}[t]{0.48\textwidth}
	  \centering
	  \includegraphics[width=\textwidth]{/home/mathisdr/Desktop/2024-2025/Stage/Jupyter_stage/Figures/Godfray_effect-of-n1n3.pdf}
	  \caption{}
	  \label{fig:a}
	\end{subfigure}
	\hfill
	\begin{subfigure}[t]{0.48\textwidth}
	  \centering
	  \includegraphics[width=\textwidth]{/home/mathisdr/Desktop/2024-2025/Stage/Jupyter_stage/Figures/Godfray_effect-of-n2m.pdf}
	  \caption{}
	  \label{fig:b}
	\end{subfigure}

\caption{Effect of maturation variability on different properties of the Godfray and Hassel 1989 model.  In all subfigures, parameters are those taken in Xu et al. 2010.  \textbf{(a)} : Local stability diagram for $(n_2,m) = 3$ and various values of $(n_1,n_3)$. \textbf{(b)} : Local stability diagram displaying two couples of conjugated eigenvalues responsible for the GC area at $ratioT \approx 1.5$.}
\label{fig:Godfray}
\end{figure}

\tabto{1 cm} Similary to what was found with the precedent system, maturation time variability of the various stage increases the non-trivial equilibrium densities (especially at low ($n_2, m$)) (Figure 3a). Maturation variability can also increase the amplitude and period of oscillations ([SUPPL FIGURE]. 
\tabto{1 cm} Xu et al. (2010) showed that GCs are not possible for ($n_2,m$) = 11 with a generation time ratio superior to 1 and that the GC area is less sensitive to (Erlang-)distributed delays of $H_2, P_1$ for a ratio around 1/2. Here, we show that : \\
\tabto{1.5 cm} 1. At ($n_2,m$) = 3, the 1/2-ratio-GC area is greatly reduced when ($n_1, n_3$) are smaller than 100--300 (Figure 3a). This means that when ratioT $\approx$ 1/2, GC can occur only if the developmental time variabilities of the "waiting" stages $H_1, H_3$ are extremely low. \\
\tabto{1.5 cm} 2. While fixing $(n_1, n_3) = 300$, GC area at ratio $\approx$ 1.5 reappear for ($n_2,m$) between 50 and 100 (Figure 3b). Together with the results of Xu et al. (2010) that showed the low effect of m compared to $n_2$, this means that for ratioT $\approx$ 1.5, 2.5 ... , GC can occur only if the developmental time variability of the vulnerable stage $H_2$ is very low. \\
-> Refaire cette petite partie dure à lire sous forme d'un tableau récapitulatif ? comment le structurer ? 























\section{Discussion (d'abord sous forme de bullet points}
Soit je le structure en mode "n fait changer les modèles dont il ne faut pas l'ignorer et l'estimer" soit l'inverse : "on a ces n là, et les modèles sont +- sensibles, donc ...". Je vais plutôt choisir la première option qui a plus de sens du pdv du storytelling. Mais j'ai tjrs un gros pb avec ça : avec les travaux de Xu et de Wearing, j'ai vraiment l'impression de retrouver naivement ce qu'ils ont déjà dit. En plus, je "spoile" déjà dans l'intro en ayant déjà introduit le concept de sensibilité.
\\
\textbf{Mais donc si je prends cet ordre je devrais surement le prendre aussi pour les résultats?}

\subsection{Understanding the effect of developmental variability}
1.Effet de n sur $p_{survie}$ (cf oral de Juin)
\\
2.Ressemblance frappante entre l'effet de l'agrégation spatiale du parasitisme (ajout de pseudointerférénce via k dans un modèle à délai fixe) et l'effet de l'agrégation temporelle du parasitisme (ajout de variab dev. via n dans un modèle sans pseudointerférence) [ici pb car il faut justifier d'où vient, comment a été construite la figure...]



\begin{figure}[!ht]
\centering	
\includegraphics[width=0.8\textwidth]{/home/mathisdr/Desktop/2024-2025/Stage/Internship_2025_LCT/model_host_parasitoid/figs/Equivalence_k_n.pdf}
\caption{Effect of model parasitism heterogeneity on the stability of the Murdoch 1987. n is the number of subcompartiments, or the shape fo the Gamma distribution of maturation time. In case of $n \rightarrow \infty$, maturation time is fixed and parasitism risk is homogeneous over all hosts. k is the term of overdispersion of parasitism which is dispersed . In case of $k \rightarrow \infty$, parasitism is perfectly homogenenous over all hosts.}
\label{fig:equivalence_k_n}
\end{figure}

\subsection{Differential sensitivity of models to developmental variability}
Gurney 1980 model ; Larval cohort competition : model built by Gurney 1980 and further justified in Gurney 1983. Equivalent Erlang-distributed model analytically studied by Beretta et al. 2016. Les résultats sont similaires : on a une évolution "non linéaire" des prop du système avec une apparition brutale d'une zone d'instabilité quand n augmente. Mais cette évolution est dépendante d'autres paramètres (ce qu'illustre très bien Beretta et al. 2016). Idem dans Wearing : changement de type de dynamique entre GC et MGC.
\tabto{1.5 cm} -> Let sensitivity (understand : to the increase of n) be the propensity of ODE models’ prediction to swiftly change as soon as maturation time variability is reduced (as long as n increases from 1).  Note that this definition of sensitivity is inverse to the one chosen by Wearing et al. In any case, they showed that sensitivity was variable and dependent on several factors (to which developmental time variability is added ; to which dynamic is studied ; to the used set of parameters and their distance to a pre-existing bifurcation).  [... et nous ce qu'on montre c'est exactement la même chose : évolution "non linéaire" du système avec n, mais ça dépend de plein de choses... donc c'est du cas par cas et faut faire gaffe : on a rien réinventé finalement on confirme juste Wearing pour moi? C'est un peu ça qui me donne des difficultés pour structurer ce rapport]
\\
Murdoch comme Gurney: exemple typique du modèle sensible (avec réf ODE). Remarquer notamment la sensib de l'effet de la fécondité sur la variab dev : important du pdv de la modélisation en biocontrôle.
\\
\\
\textit{Godfray and co-workers collected examples of cycles in tropical insects where the cycle period is roughly equal to the development time of the host, and no climatic fluctuation could account for the oscillating dynamics (Godray and Hassel 1989). They proved that phenological asynchrony can cause single-generation cycles ("GC") in populations (Godfray and Hassel 1989, Godfray et al.1994). This happens if the ratio between host and parasitoid development times leads to a gap between host and parasitoid density peaks. This will lead to high parasitism at times when the host population is at a low point, further reinforcing the trough (Murdoch 2003).  To reproduce this behavior in a "strategic" model, they developed a DDE model that indeed predicts GCs for generation time ratios equal to 0.5, 1.5, ... J'avais écrit ça mais on s'en fout un peu ici non? Enfin c'est déjà un peu introduit dans les résultats, pas besoin de s'attarder sur cette théorie-là? } \\
\tabto{1.5 cm} -> Par contre il faut pour moi dire un truc de ce style : revenir sur le côté +- "artéfactuel" de ce que j'appelle 'l'effet godfray' qui reste très dép d'un délai parfaitement fixe pour tous les compartiments considérés. A la fois la fixité de H2 et P1 : si pas fixe les GC à ratioT>1 disparaissent. Mais comme les temps de dev des parasitoides sont en général plus court que ceux des hôtes, c'est moins important/fréquent dans la nature probablement. Par contre pour ratioT autour de 0.5 : plus robuste à l'ajout de variab dev, mais pas non plus infiniment robuste. Tres dép de des stades d'attentes H1 et H3 qui ajoutent du délai dans le système. SI on les supprime (j'ai testé personnellement, ça fait remonter +++ les courbes de stabilité en U), ou si on ajout un peu de variab, on perd cette zone GC. 
\\
\\
\\
-> (transi vers la suite) Donc faut estimer cette variab sur le système qu'on étudie si [on veut pas prendre le risque de rater des propriétés majeures du système] (mal dit?).


\subsection{Variability of maturation time in laboratory and in nature}
Retour sur les résultats : figure 1a and 1b show that [à quel point répeter tout? est-ce nécessaire de refaire un summary?]
	2.effet de constraint : bcp de n envoyés à droite si tau = 0 (avec supplementary graph qui illustre). Ainsi si pas de shift on aura jamais de n en dessous de . 
	3.n de total dev > mais pas forcément >>>>>.
	4.eventuellement parler de la cphérence avec Xu : on pourra supplementary graph un clipped histogram : même en essayant de reprod oui on a un groupe à 1-10, mais on a bcp bcp + de gens à >20!
\\
Autres manières d'estimer la variab développementale :
	-Trouver des articles où le CV du temps de maturation est donné et appliqué la formule donnée par la loi Gamma (ce que fait Jules)
	-La technique de Blythe avec la décomposition tau = tau1+tau2 et le rapport tau1/tau2 : on pourrait calculer ce rapport pour nos différentes shapes de loi Gamma (et shift) ? -> En tout cas Blythe a aussi trouvé des rapports très variés (par rapport à un thresold acceptable de 5 qu'il propose via son modèle)  qui soulignent que c'est un peu du cas par cas. On est d'accord avec lui là-dessus.
\\
On a donc vu que beaucoup de prédictions données par le modèle "simple" H-P Murdoch changent brutalement dans la zone de n 1-10. Or, on voit dans le survey empirique, lorsqu'on n'autorise pas de pré-shift, qu'il n'existe pas de n autour de 1-6. Ca prouve donc que l'approche ODE est effectivement peu réaliste et apporte un réel biais dans les prédictions des modèles, à la fois quantitatives (equilibres, propriétés des oscillations) et qualitatives (types de dynamiques). 
\\
Pour autant, les dynamiques GC du modèle de Godfray suggèrent qu'on ne peut pas conclure non plus que l'approche DDE est nécessairement plus acceptable autant du pdv du réalisme que de la robustesse des prédictions des modèles à délaix fixes à l'ajout d'un tout petit peu de variab dev.

\section{Perspectives}
Survey-CV : une manière simple et rapide d'estimer le nombre de sous-compartiments à appliquer.... quand le LCT est nécessaire 
Applications en modèles non autonomes pr prédire les outbreaks : travail de Jules
Reseaux d'intéractions?







\subsection{Funding : en attente d'un meilleur template?}
The internship is funded by the Cacolac project (OFB, Ecophyto...), which objective is to integratively develop agroecological methods of pest control. It aims to better understand the complementarity or antagonism between different operational levers used to control bio-aggressors in perennial crops and to optimize their combination according to the local context and landscape.


% References (excluded from page count)
\newpage
\noPageNumber
\begin{thebibliography}{99}
\bibitem{example} Author, \emph{Title}, Journal, Year.
\end{thebibliography}





% Appendix (excluded from page count)
\newpage
\noPageNumber
\appendix
\section{Appendix}

\subsection{Murdoch 1987 system}
\subsubsection{System construction : application of the Linear Chain Trick with a non-autonomous mortality term}
Note that Hurtado et al. 2018 as well as Lou and Sun 2022 give detailed justifications about the construction of continuous stage-structured models with an arbitrary maturation distribution and their convergence to a coupled system of ODEs in the case of the Erlang distribution. However, they do not include a non-autonomous mortality term. Let us thereby justify thoroughly how the studied LCT model was built analogously to the DDE model given by Murdoch (1987).

Let us consider the general form of the models used both in the ODE and DDE framework :  
\begin{equation}
\dot{H_{1}}(t) = B(t) - M_n(t) - \delta(t, ...) H_1(t)
\end{equation}

$H_1(t)$ is the density of juvenile hosts. $B(t)$ is the recruitment term of the juveniles, ie the birth rate. $M_n(t)$ is the maturation term. In the ODE framework, it would be of the form $M_{ODE}(t)= m_H H_1(t)$. In the DDE framework, it would be of the form $M_{DDE}(t) = B(t-\tau_H)S_{DDE}(t)$ with $S_{DDE}(t) := \exp\left( -\int_{t-\tau_H}^{t} \delta(u, ...) \, du \right)$ (Gurney and Nisbet 1983). 

$\delta_H(t, ...)$ is a mortality term. \textbf{We assume that $\delta$ can depend on both the densities of host and parasitoids}. Though, for clarity, $\delta(t, ...)$ will now be written in the form $\delta(t)$.

We give to the equation the initial conditions $H_1(0) = H_1^0$ and $B(-t) = B^{-}(t)$, $\delta(-t) = \delta^{-}(t)$, for $t \geq 0$.

To account for a arbitrary residence time distribution, the maturation term can be written in an integral form :
\begin{equation}
M_n(t) = \int_0^{+\infty} B(t-s)\varphi_n(s) \exp\left( -\int_{t-s}^{t} \delta(u) \, du \right) ds
\end{equation}
where $\varphi_n$ is the probability density function (PDF) of a gamma distribution of rate $\theta_H$ and shape n. Let's recall that

\begin{align}
\varphi_1(0) &= \theta_H, & \varphi_1' &= -\theta_H \varphi_1 \\
\varphi_n(0) &= 0, & \varphi_n' &= \theta_H (\varphi_{n-1} - \varphi_n) \quad \text{for } n \geq 2
\end{align}

Let us write :
\[
S(s, t) := \exp\left( -\int_{t-s}^{t} \delta(u) \, du \right)
\]
which satisfies :
\begin{equation}
S(0, t) = 1, \quad \partial_s S(s,t) + \partial_t S(s,t) = -\delta(t) S(s,t)
\end{equation}

We can thus write :
\begin{equation}
M_n(t) = \int_0^{+\infty} B(t-s)\varphi_n(s)S(s,t) \, ds
\end{equation}

By derivation :
\[
M_n'(t) = \int_0^{+\infty} B'(t-s)\varphi_n(s)S(s,t) \, ds + \int_0^{+\infty} B(t-s)\varphi_n(s) \partial_t S(s,t) \, ds
\]

The first term, with an integration by part in s and using $S(0,t) = 1$, gives :
\begin{align*}
\int_0^{+\infty} B'(t-s)\varphi_n(s)S(s,t) \, ds &= -[B(t-s) \varphi_n(s) S(s,t) ]_{0}^{\infty} + \int_0^{+\infty} B(t-s)\frac{d}{ds}(\varphi_n(s)S(s,t)) \, ds \\
&\quad = -0 + \varphi_n(0) B(t) + \int_0^{+\infty} B(t-s)\varphi_n'(s) S(s,t) \, ds \\
&\quad + \int_0^{+\infty} B(t-s)\varphi_n(s) \partial_s S(s,t) \, ds
\end{align*}

With $(4)$, we have :
\begin{align*}
&\int_0^{+\infty} B(t-s)\varphi_n(s) \partial_s S(s,t) \, ds  + \int_0^{+\infty} B(t-s)\varphi_n(s) \partial_t S(s,t) \, ds \\
&\quad = - \delta(t) \int_0^{+\infty} B(t-s)\varphi_n(s) S(s,t)  \, ds \\
&\quad = - \delta(t) M_n(t)
\end{align*}

we finally get :
\[
M_n'(t) = \varphi_n(0) B(t) + \int_0^{+\infty} B(t-s)\varphi_n'(s) S(s,t) \, ds - \delta(t) M_n(t)
\]

In the case $n = 1$, i.e in the ODE framework, using $(3)$ : 
\begin{align*}
M_{1}'(t) &= M_{ODE}'(t)= -\theta_H B(t) + \int_0^{+\infty} B(t-s)(-\theta_H \varphi_1(s)) S(s,t) \, ds - \delta(t) M_1(t) \\
&= \theta_H B(t) - (\theta_H + \delta(t)) M_1(t)
\end{align*}

And for $n \geq 2$, using $(4)$  :
\[
M_n'(t) = \theta_H M_{n-1}(t) - (\theta_H + \delta(t)) M_n(t)
\]

Thus, \(M_n\) is solution of the system :
\begin{equation}
\begin{cases}
\dot{M}_1(t) = \theta_H B(t) - (\theta_H + \delta(t)) M_1(t) \\
\dot{M}_k(t) = \theta_H M_{k-1}(t) - (\theta_H + \delta(t)) M_k(t), \quad k \in \{2, \dots, n\}
\end{cases}
\end{equation}

The initial conditions of this system  $(M_1(0), \dots, M_N(0))$ are given by $(2)$, with $t = 0$.
Now, let us write, for \(k \in \{1, \dots, N\}\), $H_{1,k}(t) := \frac{M_k(t)}{\theta_H}$. Then : 
\begin{equation}
\begin{cases}
\dot{H_{1,1}}(t) = B(t) - (\theta_H + \delta(t)) H_{1,1}(t) \\
\dot{H_{1,1}}(t) = \theta_H H_{1, k-1}(t) - (\theta_H + \delta(t)) H_{1,k}(t), \quad k \in \{2, \dots, n\}
\end{cases}
\end{equation}

By summing all the lines, and writing $y(t) := H_{1,1}(t) + \dots + H_{1,n}(t)$, we get :
\[
\dot{y}(t) = B(t) - \theta_H H_{1,n}(t) - \delta(t) y(t) = B(t) - M_n(t) - \delta(t) y(t)
\]

Thus, $H_{1,1} + \dots + H_{1,n}$ does satisfy the initial condition $(1)$.

In particular, if $H_{1,1}(0) + \dots + H_{1,n}(0) =  H_1^0$, then $H_1(t) = y(t) = H_{1,1}(t) + \dots + H_{1,n}(t)$ and $(1)$ is equivalent to $(8)$. This is called the Linear Chain Trick.

Applying the same method on the parasitoid maturation term with Erlang rate $\theta_P$ and shape $m$, choosing a linear parasitism term of parasitism assuming a random and homogeneous search of hosts by the parasitoids (May 1978), and assuming a linear birth rate, the final LCT system can be rewritten as a system of coupled ODEs  :
\begin{equation}
\begin{dcases}
H_{1,1}'(t) = H_{1,1}(t)(-a P_2(t) - \delta_{H1} - \theta_H) + r H_2(t) \\
\vdots \\
H_{1,i}'(t) = \theta_H(H_{1,i-1}(t) - H_{1,i}(t)) - \delta_{H1} H_{1,i}(t) - aP_2H_{1,i}(t), \quad i \in [[2,n]] \\
\vdots \\
H_{2}'(t) = \theta_H H_{1,n}(t) - \delta_{H2} H_{2}(t) \\
P_{1,1}'(t) = P_{1,1}(t)(- \delta_{P1} - \theta_P) + a P_2(t) \sum_{i=1}^{n} H_{1,i}(t) \\
\vdots \\
P_{1,i}'(t) = \theta_P(P_{1,i-1}(t) - P_{1,i}(t)) - \delta_{P1} P_{1,i}(t), \quad i \in [[2,m]] \\
\vdots \\
P_{2}'(t) = \theta_P P_{1,m}(t) - \delta_{P2} P_{2}(t) \\
\end{dcases}
\end{equation}

\subsubsection{Equilibria}

Setting the right side of the model system equal to zero and solving for equilibrium yields the following non-trivial equilibrium densities :
\begin{align}
H_1^* &= \frac{\delta_{P2}}{a(f_P)^m} \nonumber \\
H_2^* &= H_1^* * \frac{\theta_{H} (\rho^{\frac{1}{n}}-1)}{\delta_{H2} (\rho-1)} \nonumber \\
P_1^* &= P_2^* * \frac{\alpha_P \delta_{P2}}{\theta_P (f_P)^{m-1}}  \\
P_2^*&= \frac{\theta_H}{a f_H} (\rho^{\frac{1}{n}}f_H -1) \nonumber
\end{align}

With :
\begin{align}
f_* &:= \frac{\theta_*}{\theta_* + \delta_*}, \quad \text{$_*$ equals H,P} \nonumber  \\
s_H^* &:= \frac{\theta_H}{\theta_H + \delta_H + aP_2^*} = \rho^{-\frac{1}{n}} \nonumber \\
\alpha_H^* &:= \sum_{i=0}^{n-1} (s_H^*)^i   \\
\alpha_P^* &:= \sum_{i=0}^{m-1} (f_P^*)^i \nonumber
\end{align}

These equilibria are strictly positive if :
\begin{equation}
(\frac{\beta}{\delta_{H2}})^{\frac{1}{n}} f_H  > 1
\end{equation}

\subsubsection{Parametrisation}
To find the mean development time, or expected lifespan of juveniles that reach maturity, the Erlang distribution properties can be used.
For clarity, let us reason for hosts in the following equations. Parasitoid terms are set in the same way.  Let us recall that the maturation time is Erlang-distributed with rate $\theta_H$ and shape n. For juveniles, the mean maturation time $\hat{\tau_{H}}$, the variance of maturation time $\sigma^2_{H}$ and the coefficient of variation of maturation time $CV_{H}$ would be : 
\begin{align}
\hat{\tau_{H}} = \frac{n}{\theta_{H}} \nonumber \\
\sigma^2_{H} = \frac{n}{(\theta_{H})^2} \\
CV_{H} = \frac{1}{\sqrt{n}} \nonumber 
\end{align}

Setting  $\hat{\tau_{H}} = 1/m_H$, the maturation term becomes :
\begin{equation}
\theta_H = n*m_H
\end{equation}

Also, setting $\rho$ the mean number of offspring yielded by a host individual during its adult lifetime, we get :
\begin{equation}
r = \rho \delta_{H2}
\end{equation}

However, in the previous lines, juvenile intrinsic mortality was not taken in account. Alternatively, a continuous-time stochastic agent-based model can be built. This intuitively accounts for the stochastic process lying behind the model assumptions.

Consider a juvenile agent that either matures or dies. Maturation time $\tau_m$ is drawn from the Erlang distribution $\varphi(\theta_H, n)$ and death time $\tau_d$ is drawn from the Exponential distribution $\mathcal{E}(\delta_H)$. If $\tau_m \leq \tau_d$ then the agent matures. Once adult, new juveniles are produced from a Poisson point process of rate $\beta$.

Thus, let $f_{surv,H1}$ be the mean proportion of surviving host individuals and (13) become :

\begin{align}
\hat{\tau_{H}} = \frac{n}{\theta_{H} + \delta_{H}} \nonumber \\
\sigma^2_{H1} = \frac{n}{(\theta_{H} + \delta_{H})^2} \\
CV_{H} = \frac{1}{\sqrt{n}} \\
f_{surv,H1} = (\frac{\theta_{H}}{\theta_{H} + \delta_{H}})^n
\end{align}

Still setting $\hat{\tau_{H}} = 1/m_H$, the maturation and intrinsic mortality terms then become :
\begin{equation}
\begin{alignedat}{2}
\theta_{H} = n * m_{H}* (f_{surv,H1})^{1/n} \\
\delta_{H} = n* m_{H}* (1-(f_{surv,H1})^{1/n})
\end{alignedat}
\end{equation}

We notice that the existence of a positive non-trivial equilibrium can be rewritten as :
$$\rho^{\frac{1}{n}}f_H = (\rho f_{surv,H1})^{\frac{1}{n}} > 1 \qquad \impliedby \qquad \rho f_{surv,H1} > 1$$

\subsubsection{Convergence to ODE and DDE models}
If $n,m = (1,1)$, the Erlang distribution is equal to an exponential distribution and the system is equivalent to its simple ODE system counterpart. Then, $f_*$, $s_H$, $\alpha_*$ are all equal to 1 and equilibrium densities are those found in the simple ODE system. \\ \\

If $n,m\to\infty$, the Erlang distribution of maturation time converges to a dirac distribution and the system is equivalent to its DDE counterpart. Then, equibrium densities become : 
\begin{align}
H_1^* & \xrightarrow[n,m\to\infty]{} \frac{\delta_{P2}}{a\gamma} \nonumber \\
H_2^* & \xrightarrow[n,m\to\infty]{} \frac{ln(\rho)}{\rho-1} \frac{\delta_{P2} m_H}{a \gamma \delta_{H2}} \nonumber \\
P_1^* & \xrightarrow[n,m\to\infty]{} P_2^{* \infty} \frac{\gamma -1}{\gamma} \frac{\delta_{P2}}{m_P ln(\gamma)} = \frac{(\gamma-1)m_H(ln(\rho) - \frac{\delta_{H1}}{m_H})}{\gamma  a  m_P   ln(\gamma)} \\
P_2^*&  \xrightarrow[n,m\to\infty]{} \frac{m_H}{a}(ln(\rho) - \frac{\delta_{H1}}{m_H}) \nonumber  \\
\text{With} \quad \gamma & := e^{-\frac{\delta_{P1}}{m_P}} \sim_{m \to \infty} (f_P)^m \nonumber
\end{align}


Expectedly, they are strictly equivalent to equilibrium densities given in Murdoch 1987.

\subsection{Godfray 1989 system}

Setting the right side of the model system equal to zero and solving for equilibrium yields the following non-trivial equilibrium densities:
 \begin{align}
  H_1^* &= H_2^* * \frac{\theta_{H2}+\delta_{H2}+\Pi^*}{\theta_{H1}} \frac{\alpha_{H1}}{\alpha_{H2}} (f_{H1})^{1-n_1} \nonumber \\
  H_2^* &= P_2^* * \frac{\delta_{P2}}{(f_{P1})^{m} \Pi^*} \nonumber \\
  H_3^* &= H_4^* * \frac{\delta_{H4} \alpha_{H3}}{\theta_{H3}(f_{H3})^{n_3-1}}  \nonumber \\
  H_4^* &= H_1^* * \frac{\theta_{H1}+\delta_{H1}}{r \alpha_{H1}}  \\
  P_1^* &= H_2^* * \frac{\Pi^*}{\delta_{P1}}(1-(f_{P1})^{m}) \nonumber \\
  P_2^* &= \frac{k}{a}exp(\frac{\Pi*}{k}-1) \nonumber 
\end{align}

With :
 \begin{align}
  f_X^* &:= \frac{\theta_X}{\theta_X + \delta_X}, \quad \text{$X$ equals H1,H3,P1} \nonumber \\
  s_{H2}^* &:= \frac{\theta_{H2}}{\theta_{H2} + \delta_{H2} + \Pi^*} = [\frac{\delta_{H4}}{r}(f_{H1})^{-n1}(f_{H3})^{-n_3}]^{\frac{1}{n_2}}\nonumber \\
  \alpha_X^* &:= \sum_{i=0}^{n_X-1} (f_X^*)^i,  \quad \text{$X$ equals H1, H3,P1} \\
  \alpha_{H2}^* &:= \sum_{i=0}^{n_2-1} (s_{H2}^*)^i \nonumber \\
  \Pi^* &:= k*ln(1+\frac{aP_2^*}{k}) = \frac{\theta_{H2}}{f_{H2}} ((\frac{r}{\delta_{H4}}(f_{H1})^{n1}(f_{H3})^{n3})^{\frac{1}{n2}}f_{H2} -1)\nonumber 
 \end{align}


\subsection{Supplementary figures}
\renewcommand{\figurename}{SFig}
\setcounter{figure}{0}


\begin{figure}[!ht]
\centering	
\includegraphics[width=0.8\textwidth]{/home/mathisdr/Desktop/2024-2025/Stage/Internship_2025_LCT/model_host_parasitoid/figs/Murdoch_bifurcation_diagram_Other-par.pdf}
\caption{Bifurcation diagram of the system with parameters inspired from those in Murdoch 1987 : $\delta_{H1} = 0, \delta_{P1} = 1, \delta_{P2} = 8, \rho = 33, m_P = 2,  T_{H2} = 1/\delta_{H2} = 2, T_{H1} = 1/m_H = 1$.}
\label{fig:bifurc_suppl}
\end{figure}



\end{document}
